\documentclass[11pt,a4paper]{article}
\usepackage[margin=2.2cm]{geometry}
\usepackage{amsmath,amssymb}
\usepackage{graphicx}
\usepackage{booktabs}
\usepackage{xcolor}
\usepackage[colorlinks=true,linkcolor=blue!50!black,urlcolor=blue!50!black]{hyperref}
\usepackage{enumitem}
\usepackage{fontspec}
\setmonofont{Menlo}[Scale=0.82]
\setlist{itemsep=2pt,topsep=4pt}
\setlength{\emergencystretch}{3em}

\newcommand{\code}[1]{\texttt{#1}}

\title{Immersed interface method in StagFDTools\\[2pt]
\large Notes on steps 1--8: from interface geometry to second-order incompressible and
compressible inclusion tests, in Landry's z-form, in a StagFDTools stress form, and matrix free}
\author{StagFDTools, branch \texttt{iim-discontinuous-viscosity}}
\date{\today}

\begin{document}
\maketitle

\noindent This note documents the first steps of the implementation plan
(\texttt{docs/IIM\_plan.md}) for porting the immersed interface method (IIM) of
W.~Landry (\url{https://gitlab.com/wlandry/discontinuous_stokes}) to StagFDTools.
In that method, the Stokes equations are discretised on the standard MAC grid in terms
of $z=\eta\mathbf v$ and $p$. Stencils that cross the viscosity interface $\Gamma$ receive
correction terms built from the jumps $[z]$, $[\nabla z]$ and $[p]$.
\begin{itemize}
  \item \textbf{Step~1} (Sections~\ref{sec:s1impl}--\ref{sec:s1res}) provides the geometry
    these corrections need: where $\Gamma$ crosses the grid, the normal there, and which
    equations are affected.
  \item \textbf{Step~2} (Section~\ref{sec:s2}) provides the interface velocity
    $\mathbf v_\Gamma$ and its tangential derivative $\partial_t\mathbf v_\Gamma$, on which
    all the jumps depend.
  \item \textbf{Step~3} (Section~\ref{sec:s3}) turns them into the jumps $[z]$, $[\nabla z]$
    and $[p]$.
  \item \textbf{Step~4} (Section~\ref{sec:s4}) assembles the corrected operator, solves the
    viscous inclusion test, and compares the result with standard FD.
  \item \textbf{Step~5} (Section~\ref{sec:s5}) adds quadratic fits for the interface traces
    and second-order corrections.
  \item \textbf{Step~6} (Section~\ref{sec:s6}) extends the method to compressible viscous flow
    and tests it on circular and elliptical inclusions.
  \item \textbf{Step~7} (Section~\ref{sec:s7}) reformulates the IIM with the kernels of the
    StagFDTools \code{Stokes} module (velocity interpolation, stresses at centroids and
    vertices), as a basis for anisotropic and non-linear rheologies, and compares it with the
    z-form.
  \item \textbf{Step~8} (Section~\ref{sec:s8}) solves the stress form matrix free with dynamic
    relaxation, in the vectorised style of the standalone PHDR codes.
\end{itemize}
The code is in \code{src/StokesIIM/}, exported by the submodule \code{StagFDTools.StokesIIM}.

%---------------------------------------------------------------
\section{Discretisation}

Figure~\ref{fig:mesh} summarises the discretisation. There are no unknowns beyond the
standard ones: $z_x$ on the $x$-faces, $z_y$ on the $y$-faces and $p$ at the centroids.

An equation is \emph{irregular} if its stencil has members on both sides of $\Gamma$. For each
such equation, the jumps are evaluated where $\Gamma$ crosses the segment between stencil
members, and they depend on two interface quantities:
\begin{itemize}
  \item the interface velocity $\mathbf v_\Gamma$;
  \item its tangential derivative $\partial_t\mathbf v_\Gamma$.
\end{itemize}
These are not solved for. They are eliminated by one-sided linear fits on each side, combined
with a viscosity-weighted average:
\[
\mathbf v_\Gamma=\frac{\eta_-\mathbf v_-+\eta_+\mathbf v_+}{\eta_-+\eta_+}.
\]

\begin{figure}[h]
  \centering
  \includegraphics[width=\textwidth]{figures/iim_mesh.png}
  \caption{IIM discretisation on the MAC grid (script \code{docs/figures/iim\_mesh.jl}).
  \textbf{(a)}~Unknowns coloured by phase (orange: inclusion $\eta_-$, blue: matrix $\eta_+$).
  Outlined nodes are irregular equations. Gold diamonds are the crossings of $\Gamma$ with the
  half-grid lines, where jumps are evaluated; they are not unknowns.
  \textbf{(b)}~One irregular $x$-momentum equation. The star and hexagon are the interface
  points it uses for $\partial_{xx}z_x,\partial_{yy}z_x$ and for the cut corner of
  $\partial_{xy}z_y$. Crosses mark the fit centres $x_\Gamma\pm L\mathbf n$. Triangles show the
  3-node one-sided fits (solid: $z_x$, dashed: $z_y$) that give $\mathbf v_\Gamma$ and
  $\partial_t\mathbf v_\Gamma$.}
  \label{fig:mesh}
\end{figure}

%---------------------------------------------------------------
\section{Step 1: implementation}\label{sec:s1impl}

\paragraph{Convention.} $\Gamma=\{\varphi=0\}$, with $\varphi<0$ inside the inclusion and
$\varphi>0$ in the matrix. The function \code{side($\varphi$)} (line~14) returns $-1$ or $+1$.
It never returns 0, so every node belongs to exactly one phase.

\begin{enumerate}[leftmargin=*]
\item \textbf{Sampling the level set} -- \code{LevelSet(φfun, X)} (line~22).
  It evaluates $\varphi$ on the extended grids of \code{GenerateGrid}, ghost nodes included,
  and returns the NamedTuple \code{(x, y, c, v)} for the $V_x$, $V_y$, centroid and vertex
  arrays, following the StagFDTools convention.
  \begin{itemize}
    \item Nodes with $|\varphi|<10^{-10}\Delta$ lie on $\Gamma$; they are moved to the $+$ side.
    \item Without this, a node exactly on $\Gamma$ (e.g.\ the vertex $(0,2)$ for $r_c=2$,
      $\Delta=0.5$) gets $\varphi=-7\cdot10^{-16}$ by round-off. That creates an inside sliver
      of zero width, which stencils cross twice.
  \end{itemize}

\item \textbf{Fine lattice} -- \code{FineLattice(φ, X, Δ)} (line~38).
  The four staggered arrays are merged into one array $\Phi_{IJ}$ with spacing $\Delta/2$:
  \begin{center}
  \begin{tabular}{@{}ll@{}}
    \toprule
    $(I,J)$ parity & node \\
    \midrule
    odd/odd & vertex \\
    even/even & centroid \\
    odd/even & $V_x$ \\
    even/odd & $V_y$ \\
    \bottomrule
  \end{tabular}
  \end{center}
  Every interface point used by an IIM stencil lies on an edge of this lattice:
  \begin{itemize}
    \item $V_x$--$V_x$ and $V_x$--$p$ segments;
    \item the box edges of the mixed derivative;
    \item the centroid--face segments of continuity.
  \end{itemize}
  The crossings are therefore computed once, on the lattice edges, rather than separately for
  each equation.

\item \textbf{Crossings} -- \code{InterfaceCrossings(fl)} (line~123).
  On every lattice edge where \code{side} changes:
  \begin{itemize}
    \item build the cubic Lagrange polynomial through the 4 samples $k-1,\dots,k+2$ along the
      grid line (functions \code{crossing\_fraction}, line~79, and \code{lagrange}, line~61);
    \item near the ends of the lattice, where only 3 samples exist, use the quadratic instead;
    \item locate the root on $[0,1]$ by bisection. This always succeeds, because the
      polynomial takes the endpoint values, which have opposite signs.
  \end{itemize}
  The output is:
  \begin{itemize}
    \item \code{tx}, \code{ty}: the fraction along each horizontal or vertical edge
      (\code{NaN} if the edge is not crossed);
    \item \code{points}: a list of $(x_\Gamma,\mathbf n,\kappa,I,J,\text{dir})$, one per crossing (the curvature $\kappa$ is added in step~5).
  \end{itemize}

\item \textbf{Normals} -- function \code{∇Φ} (line~93).
  \begin{itemize}
    \item At both ends of the edge, $\nabla\varphi$ is approximated by central differences
      (lattice spacing $\Delta/2$, so a stencil width of $\Delta$).
    \item These are interpolated linearly to $x_\Gamma$ and normalised.
    \item $\mathbf n$ points towards $\varphi>0$, i.e.\ out of the inclusion.
  \end{itemize}

\item \textbf{Irregular equations} -- \code{IrregularNodes(φ, nc)} (line~161).
  An equation is flagged when any member of its standard stencil lies on the other side of
  $\Gamma$. The indexing matches \code{Stokes.jl}:
  \begin{center}
  \begin{tabular}{@{}ll@{}}
    \toprule
    equation & stencil members checked \\
    \midrule
    $V_x[i,j]$ & $V_x[i\pm1,j]$, $V_x[i,j\pm1]$, $P[i-1{:}i,\,j-1]$, $V_y[i{:}i+1,\,j-1{:}j]$ \\
    $V_y[i,j]$ & $V_y[i\pm1,j]$, $V_y[i,j\pm1]$, $P[i-1,\,j-1{:}j]$, $V_x[i-1{:}i,\,j{:}j+1]$ \\
    $P[i,j]$ & $V_x[i{:}i+1,\,j+1]$, $V_y[i+1,\,j{:}j+1]$ \\
    \bottomrule
  \end{tabular}
  \end{center}

\item \textbf{Thin features} -- \code{MultipleCrossings(Γ)} (line~193).
  Counts the lattice segments of length $\Delta$ that $\Gamma$ crosses twice. These are
  features thinner than about $\Delta$, which the first-order IIM does not handle (Landry's
  ``anticorners''). The count is $0$ at every resolution of the inclusion test.
\end{enumerate}

\noindent Unit tests are in \code{test/test\_iim\_geometry.jl}, which is included in
\code{runtests.jl}. They check that:
\begin{itemize}
  \item the fine lattice reproduces $\varphi$ at every node;
  \item crossings and normals are exact (to $10^{-12}$) for a straight inclined interface;
  \item the irregular columns for a vertical interface match a hand check;
  \item nodes on $\Gamma$ are moved to the $+$ side.
\end{itemize}
The example is \code{examples/\_Stokes\_IIM/IIM\_Step1\_Geometry.jl}.

%---------------------------------------------------------------
\section{Step 1: results for the viscous inclusion geometry}\label{sec:s1res}

The test geometry is that of Schmid \& Podladchikov (2003): a circular inclusion of radius
$r_c=2$ centred in a square domain of side $L=10$. The grid spacing is
$h=\Delta x=\Delta y=L/n_c$, and two level-set functions with the same zero contour are used:
\[
\varphi_d=r-r_c\quad\text{(signed distance)},\qquad
\varphi_t=\tanh(r-r_c)\quad\text{(not a distance function)}.
\]

\begin{figure}[h]
  \centering
  \includegraphics[width=\textwidth]{figures/IIM_Step1_Geometry.png}
  \caption{Geometry for the inclusion test.
  \textbf{Left:} $n_c=20$, $\varphi=\varphi_d$. Circles are $P$, right-pointing triangles $V_x$,
  upward triangles $V_y$; colours give the phase. Outlined nodes are the irregular equations.
  Gold diamonds are the 60 crossings of $\Gamma$ with the fine lattice.
  \textbf{Middle:} zoom on the first quadrant with the computed normals $\mathbf n$.
  \textbf{Right:} maximum error of crossing positions and normals against $h$, for $\varphi_d$
  and $\varphi_t$, with $O(h^2)$ and $O(h^4)$ reference slopes.}
  \label{fig:step1}
\end{figure}

\paragraph{Left and middle panels.}
\begin{itemize}
  \item The irregular equations form a band one stencil wide on each side of $\Gamma$: 48
    $V_x$, 48 $V_y$ and 24 $P$ equations at $n_c=20$.
  \item Only $V_x$ and $V_y$ nodes up to one cell away from $\Gamma$ are flagged, because the
    momentum stencils reach $\pm h$. Only cells cut by $\Gamma$ between their centroid and a
    face are flagged for continuity.
  \item The normals are perpendicular to the circle and point outward.
  \item Some crossings come in close pairs near grid nodes. These are crossings of the
    horizontal and the vertical lattice line through the same region, which is expected.
\end{itemize}

\subsection*{Discussion of the convergence (right panel)}

The errors are measured over all crossings inside the physical domain:
\[
e_x=\max_\Gamma\bigl|\,|x_\Gamma|-r_c\bigr|,\qquad
e_n=\max_\Gamma\bigl|\mathbf n-x_\Gamma/|x_\Gamma|\bigr|.
\]
Table~\ref{tab:conv} lists the values and the local orders $\log_2(e_{2h}/e_h)$.

\begin{table}[h]
\centering
\small
\begin{tabular}{@{}rrrr rr rr rr rr@{}}
\toprule
& & & & \multicolumn{4}{c}{$\varphi_d=r-r_c$} & \multicolumn{4}{c}{$\varphi_t=\tanh(r-r_c)$}\\
\cmidrule(lr){5-8}\cmidrule(l){9-12}
$n_c$ & $h$ & points & irr.\ $V_x/V_y/P$ & $e_x$ & ord. & $e_n$ & ord. & $e_x$ & ord. & $e_n$ & ord.\\
\midrule
 10 & 1      &   28 & 24/24/12    & 3.9e-4  & --   & 9.3e-3 & --   & 1.5e-3 & --   & 1.5e-2 & --   \\
 20 & 0.5    &   60 & 48/48/24    & 1.9e-5  & 4.39 & 2.4e-3 & 1.94 & 7.6e-5 & 4.28 & 4.4e-3 & 1.74 \\
 40 & 0.25   &  124 & 88/88/44    & 1.7e-6  & 3.44 & 7.9e-4 & 1.63 & 6.8e-6 & 3.48 & 1.2e-3 & 1.93 \\
 80 & 0.125  &  252 & 184/184/92  & 1.0e-7  & 4.08 & 2.1e-4 & 1.89 & 3.8e-7 & 4.18 & 2.8e-4 & 2.06 \\
160 & 0.0625 &  508 & 360/360/180 & 6.6e-9  & 3.94 & 4.9e-5 & 2.13 & 2.7e-8 & 3.79 & 7.3e-5 & 1.91 \\
320 & 0.03125& 1020 & 728/728/364 & 4.3e-10 & 3.95 & 1.3e-5 & 1.90 & 1.7e-9 & 4.00 & 1.8e-5 & 2.00 \\
\bottomrule
\end{tabular}
\caption{Convergence of interface positions and normals for the inclusion ($r_c=2$, $L=10$).
\emph{points}: number of crossings; \emph{irr.}: number of irregular equations.
No multiple crossings occur at any resolution.}
\label{tab:conv}
\end{table}

\paragraph{Positions converge at fourth order.}
\begin{itemize}
  \item Along a grid line, the crossing is the root of the cubic interpolant through samples
    spaced $\Delta/2$. The interpolation error is
    \[
    \frac{\varphi^{(4)}}{4!}\prod_k(s-s_k)=O(h^4),
    \]
    and bisection adds nothing (60 halvings).
  \item The error that matters is the distance of $x_\Gamma$ from $\Gamma$. Shifting the root
    along the line by $\delta s=\delta\varphi/|\partial_s\varphi|$ moves the point off
    $\Gamma$ by
    \[
    \delta s\,|\mathbf n\cdot\mathbf e_s|=\delta\varphi/|\nabla\varphi|.
    \]
  \item The measured error $e_x$, a distance normal to $\Gamma$, is therefore \emph{not}
    amplified at grazing intersections,
    where $\Gamma$ is nearly parallel to the grid line. Only the tangential position of the
    point is less accurate there, and the point still lies on $\Gamma$ to $O(h^4)$.
  \item The local orders oscillate between 3.4 and 4.4 at coarse resolution. As $h$ changes,
    the crossings move relative to the samples, and the factor $\prod_k(s-s_k)$ depends on
    where in the stencil the root lies. The maximum is taken over a different set of points at
    each resolution. The mean slope is 4.
\end{itemize}

\paragraph{Normals converge at second order.}
Two approximations each contribute an $O(h^2)$ error:
\begin{itemize}
  \item the central difference over a width $\Delta$, with truncation error
    $(\Delta^2/6)\,\partial^3\varphi$;
  \item the linear interpolation of $\nabla\varphi$ along an edge of length $\Delta/2$.
\end{itemize}
Even for the exact distance function, $\partial^3\varphi\neq0$ because of the curvature
$\kappa=1/r_c$. Hence $e_n\gg e_x$ in absolute terms ($10^{-5}$ against $10^{-10}$ at
$n_c=320$). The order is nevertheless clean (1.6--2.1).

\paragraph{A distance function is not required.}
$\varphi_t$ gives the same orders as $\varphi_d$, with constants larger by a factor of about
1.5--4. Both functions have $|\nabla\varphi|=1$ on $\Gamma$, so the larger constants come from
the larger higher derivatives of $\tanh$. This matters for later steps: a level set reconstructed from
markers will not be an exact distance function, and no redistancing is needed for the geometry
to converge.

\paragraph{Consequences for the IIM.}
\begin{itemize}
  \item In the first-order IIM, the corrections have the form
    $C\sim([z]+[\nabla z]\cdot\boldsymbol\delta)/h^2$ and $C_p=\pm[p]/h$, where
    $\boldsymbol\delta$ is the distance from $x_\Gamma$ to a grid node.
  \item \emph{Position error.} An error $e_x=O(h^4)$ in $\boldsymbol\delta$ perturbs $C$ by
    $O(h^2)$, which is negligible.
  \item \emph{Normal error.} An error $e_n=O(h^2)$ rotates the $(n,t)$ frame in which the jumps
    are written. This perturbs $[\nabla z]$ and $[p]=-2[\eta]\,\partial_t v_t$ by
    $O(h^2)[\eta]$, which contributes $O(h)$ to the truncation error of the irregular
    equations.
  \item The one-sided linear fits give $\mathbf v_\Gamma$ to $O(h^2)$ and
    $\partial_t\mathbf v_\Gamma$ to $O(h)$. Once divided by $h^2$ (and multiplied by
    $\boldsymbol\delta\sim h$ for the derivative term), this leaves an $O(1)$ truncation error at
    the irregular equations. That is the reason the scheme is only first order overall.
  \item The geometry contributes at most $O(h)$, one order below the interpolation error, so it
    is not expected to limit the accuracy of the first-order scheme. The second-order variant
    will also need the curvature and the second derivatives of $\varphi$. These can be obtained
    by central differences on the same lattice.
  \item The irregular equations scale linearly with $n_c$ (about $2.3\,n_c$ for each
    velocity component), so the extra cost of the corrections is $O(n_c)$ against $O(n_c^2)$
    unknowns.
\end{itemize}

%---------------------------------------------------------------
\section{Step 2: interface velocity and tangential derivative}\label{sec:s2}

\paragraph{Method.} At each crossing $x_\Gamma$ (normal $\mathbf n$, tangent
$\mathbf t=(-n_y,n_x)$), and for each velocity component:
\begin{itemize}
  \item On each side $s=\pm1$, the 3 nodes of that component's grid closest to the fit centre
    $x_\Gamma+sL\mathbf n$ define a linear interpolant of $z$. Here
    $L=1/\max(|n_x|/\Delta x,|n_y|/\Delta y)$ is one cell along $\mathbf n$, as in Landry's
    \code{compute\_dv\_dtt.cxx}.
  \item The interpolant's value and gradient at $x_\Gamma$ give $z_s(x_\Gamma)$ and $\partial_t z_s$.
  \item $\eta$ is constant on each side and $v_s=z_s/\eta_s$, so Landry's viscosity-weighted
    average reduces to
    \[
    \mathbf v_\Gamma=\frac{z_-(x_\Gamma)+z_+(x_\Gamma)}{\eta_-+\eta_+},\qquad
    \partial_t\mathbf v_\Gamma=\frac{\partial_t z_-+\partial_t z_+}{\eta_-+\eta_+}.
    \]
  \item This is exact when $\mathbf v$ is linear on each side and continuous across $\Gamma$
    (then $z_\pm=\eta_\pm\mathbf v$).
  \item Both quantities are linear in the 6 nodal values of $z$. They are therefore stored as
    weights, which the assembly of step~4 will turn into matrix coefficients.
  \item \emph{Deviation from Landry:} he takes $\partial_t z_\pm$ from separate nearest
    one-sided differences. Here it comes from the gradient of the same fit. Both are $O(h)$.
\end{itemize}

\paragraph{Implementation} (\code{src/StokesIIM/Trace.jl}).
\begin{enumerate}[leftmargin=*]
\item \code{OneSidedFit} (line~36).
  \begin{itemize}
    \item Collect the candidate nodes within $\pm4$ cells of the fit centre that lie on side
      $s$, and sort them by distance.
    \item Take the 2 closest, then the next one that is not collinear with them
      ($|\det[1,\,\delta x,\,\delta y]|>1/2$ in cell units). Unlike Landry's row/column test,
      this also rejects diagonal alignments.
    \item Return the barycentric weights $\lambda_k$ of $x_\Gamma$ and the gradient weights
      $\mathbf g_k$, both from the inverse of the $3\times3$ matrix (\code{StaticArrays}).
  \end{itemize}
\item \code{InterfaceTrace} (line~88).
  \begin{itemize}
    \item Run the two one-sided fits.
    \item Return the 6 node indices and the weights $v_k=\lambda_k/(\eta_-+\eta_+)$ and
      $t_k=\mathbf g_k\cdot\mathbf t/(\eta_-+\eta_+)$, so that
      $v_\Gamma=\sum_k v_kz_k$ and $\partial_tv_\Gamma=\sum_k t_kz_k$.
  \end{itemize}
\item \code{InterfaceTraces} (line~136). Applies \code{InterfaceTrace} to both components at
  every crossing inside the domain. Candidates are the $V_x$ and $V_y$ nodes of the domain,
  boundary nodes included (the ranges of \code{Ranges}).
\item \code{apply\_trace} (line~150). Evaluates the weights on an array $z$.
\end{enumerate}

\noindent Tests are in \code{test/test\_iim\_trace.jl}. They check that:
\begin{itemize}
  \item the weights are exact (to $10^{-10}$) for a velocity field with a kink across an
    inclined $\Gamma$;
  \item $\sum_k v_k\eta_k=1$ and $\sum_k t_k\eta_k=0$;
  \item each fit uses only nodes of its own side.
\end{itemize}
The example is \code{examples/\_Stokes\_IIM/IIM\_Step2\_Trace.jl}. It samples
$z=\eta\mathbf v$ from the exact solution and compares the reconstructed values with the exact
ones:
\begin{itemize}
  \item $\mathbf v(x_\Gamma)$ for the value;
  \item $\mathbf L_{inc}\mathbf t$ for the derivative, where $\mathbf L_{inc}$ is the uniform
    velocity gradient of the inclusion ($\partial_t\mathbf v$ is continuous across $\Gamma$).
\end{itemize}

\begin{figure}[h]
  \centering
  \includegraphics[width=\textwidth]{figures/IIM_Step2_Trace.png}
  \caption{Interface velocity for the inclusion test ($r_c=2$, $L=10$, $\eta_{mat}=1$).
  \textbf{Left:} the fits selected by the code at three crossings ($n_c=20$,
  $\eta_{inc}/\eta_{mat}=10^3$). Stars are the crossings $x_\Gamma$ and crosses the fit centres
  $x_\Gamma\pm L\mathbf n$. Triangles are the 3-node fits on the inclusion side (orange) and the
  matrix side (blue), solid for $V_x$ and dashed for $V_y$.
  \textbf{Middle, right:} maximum error of $v_\Gamma$ and $\partial_t v_\Gamma$ over all
  crossings and both components against $h$, for three viscosity contrasts.}
  \label{fig:step2}
\end{figure}

\begin{table}[h]
\centering
\small
\begin{tabular}{@{}rr rr rr rr@{}}
\toprule
& & \multicolumn{4}{c}{$\eta_{inc}/\eta_{mat}=10^{-3}$} & \multicolumn{2}{c}{$10^{6}$}\\
\cmidrule(lr){3-6}\cmidrule(l){7-8}
$n_c$ & $h$ & $e_v$ & ord. & $e_t$ & ord. & $e_v$ & ord.\\
\midrule
  20 & 0.5     & 5.4e-1 & --   & 8.2e-1 & --   & 5.6e-5  & --   \\
  40 & 0.25    & 1.9e-1 & 1.52 & 6.2e-1 & 0.41 & 5.1e-6  & 3.47 \\
  80 & 0.125   & 6.9e-2 & 1.43 & 4.0e-1 & 0.64 & 2.6e-7  & 4.28 \\
 160 & 0.0625  & 2.2e-2 & 1.62 & 2.6e-1 & 0.62 & 2.1e-8  & 3.66 \\
 320 & 0.03125 & 6.2e-3 & 1.85 & 1.6e-1 & 0.69 & 6.2e-9  & 1.73 \\
 640 & 0.0156  & 1.7e-3 & 1.86 & 9.1e-2 & 0.81 & 1.7e-9  & 1.90 \\
1280 & 0.0078  & 4.5e-4 & 1.93 & 4.7e-2 & 0.97 & 4.5e-10 & 1.90 \\
\bottomrule
\end{tabular}
\caption{Errors of $v_\Gamma$ ($e_v$) and $\partial_tv_\Gamma$ ($e_t$). For the contrast $10^3$,
all values equal those for $10^{-3}$ times $10^{-3}$, with identical orders. For $10^6$, $e_t$
equals the $10^{-3}$ values times $10^{-6}$.}
\label{tab:step2}
\end{table}

\paragraph{Discussion.}
\begin{itemize}
  \item \emph{Orders.} $v_\Gamma$ converges at $O(h^2)$ and $\partial_t v_\Gamma$ at $O(h)$,
    as expected for linear fits. This holds only asymptotically: the derivative reaches order~1
    only around $n_c=1280$, about 250 cells per radius. Below that, the observed order is
    0.4--0.8, and the error behaves as $e_t\approx h\,|\nabla^2\mathbf v|$ on the matrix side,
    where $|\nabla^2\mathbf v|\approx12/r_c$ next to $\Gamma$.
  \item \emph{Dependence on contrast.} Inside the inclusion the exact velocity is linear, so the
    inclusion-side fit is exact. The whole error therefore comes from the matrix-side fit,
    multiplied by its weight $\eta_{mat}/(\eta_{inc}+\eta_{mat})$. This explains the factors of
    $10^{-3}$ and $10^{-6}$ between the curves.
  \item \emph{What the IIM uses.} The quantity entering the corrections is the jump
    $[z]=\Delta\eta\,v_\Gamma$, with error $|\Delta\eta|\,e_v$. This error is the same for the
    contrasts $10^{-3}$ and $10^3$ (0.54 at $n_c=20$). For $10^6$ it is the same from
    $n_c\approx320$ on ($6.2\cdot10^{-3}$), once the artefact discussed below has decayed. The
    interface treatment therefore does not degrade at high contrast.
  \item \emph{Coarse $10^6$ case.} Here $e_v$ is larger than the contrast scaling predicts and
    decays at about $O(h^4)$. This is the position error of step~1, not an interpolation error:
    $x_\Gamma$ lies $e_x=O(h^4)$ off $\Gamma$, and the exact velocity has a kink there, so the
    reference value changes by $[\partial_n v]\,e_x$. It becomes visible only once the fit
    error has been multiplied by $10^{-6}$.
  \item \emph{Consequence.} The tangential derivative is the least accurate ingredient. It sets
    $[\nabla z]$ and $[p]=-2[\eta]\partial_tv_t$, and therefore the pressure jump. If step~5
    shows that it limits the pressure accuracy, the remedy is a 6-point quadratic fit (Landry's
    second-order version) or a least-squares fit over more nodes.
\end{itemize}

%---------------------------------------------------------------
\section{Step 3: jump conditions}\label{sec:s3}

The function \code{InterfaceJumps(Δη, vΓ, ∂tvΓ, n)}, in \code{src/StokesIIM/Jumps.jl}
(line~19), is a port of Landry's \code{compute\_jumps.cxx}. It uses the following conventions:
\begin{itemize}
  \item $[q]=q_+-q_-$ and $\Delta\eta=\eta_+-\eta_-$;
  \item $\mathbf n$ points towards $\varphi>0$ and $\mathbf t=(-n_y,n_x)$;
  \item $[\nabla z]_{ab}=[\partial_b z_a]$.
\end{itemize}
It assumes that $\mathbf v$ is continuous, that incompressibility holds on both sides, and that
the traction is continuous (no surface tension). Under these assumptions:
\[
[z]=\Delta\eta\,\mathbf v_\Gamma,\quad
[\partial_t z]=\Delta\eta\,\partial_t\mathbf v_\Gamma,\quad
[\partial_n z_n]=-\Delta\eta\,\partial_t v_t,\quad
[\partial_n z_t]=-\Delta\eta\,\partial_t v_n,\quad
[p]=-2\Delta\eta\,\partial_t v_t.
\]
The gradient jumps are built in the $(n,t)$ frame and rotated to $(x,y)$. All jumps are linear
in $(\mathbf v_\Gamma,\partial_t\mathbf v_\Gamma)$.

\paragraph{Test} (\code{test/test\_iim\_jumps.jl}). At 37 points on the circle, the jumps are
computed from the exact $\mathbf v_\Gamma$ and $\partial_t\mathbf v_\Gamma$ of the Schmid
solution. They are compared with the jumps of that solution measured at $r_c(1\pm10^{-7})$, for
$\eta_{inc}/\eta_{mat}=10^{-3}$, 7 and $10^4$, with and without simple shear. The errors
($10^{-6}$ to $10^{-10}$, scaled by $|\Delta\eta|$) are at the level of the $10^{-7}$ offset.
This checks the physics and the sign conventions without any discretisation error.

%---------------------------------------------------------------
\section{Step 4: corrected operator and the viscous inclusion test}\label{sec:s4}

\paragraph{Operator} (\code{src/StokesIIM/Assembly.jl}). The interior equations use the
standard MAC stencils in $z$-form:
\[
2\partial_{xx}z_x+\partial_{yy}z_x+\partial_{xy}z_y-\partial_xp=0,\quad
\partial_{xx}z_y+2\partial_{yy}z_y+\partial_{xy}z_x-\partial_yp=0,\quad
\partial_xz_x+\partial_yz_y=0.
\]
All of Landry's correction cases (second derivatives, cut sides and corners of the mixed
derivative, pressure gradient, continuity) reduce to one rule. In the equation of a node on side
$s$, every stencil member $k$ across $\Gamma$ is replaced by the extension of the $s$-side field
to $x_k$:
\[
q_k\;\to\;q_k+s\bigl([q]+[\nabla q]\cdot(x_k-x_\Gamma)\bigr)
\qquad(\text{only }[p]\text{ for pressure members}).
\]
Here $x_\Gamma$ is the crossing on the fine-lattice path between the node and $k$ that lies
closest to their midpoint.

\paragraph{Implementation.}
\begin{enumerate}[leftmargin=*]
\item \code{JumpFunctionals} (line~35) expresses, for every crossing, the jumps as linear
  functions of the 12 cloud values of its two traces ($6\,z_x+6\,z_y$, from step~2). It does
  this by applying \code{InterfaceJumps} to the columns of the trace weights.
\item \code{pair\_crossing} (line~82) selects $x_\Gamma$ for a (node, member) pair, through a
  lookup table of crossings indexed by fine-lattice edge.
\item \code{AssembleIIM} (line~110) assembles the system:
  \begin{itemize}
    \item The unknowns are numbered with \code{set\_boundaries\_template!(:all\_Dirichlet)} and
      \code{Numbering!} from the \code{Stokes} module.
    \item Each equation adds its standard coefficients, plus the correction coefficients on the
      12 cloud nodes of the selected crossing.
    \item Non-numbered (boundary and ghost) nodes go to the right-hand side.
    \item A Lagrange multiplier ($\sum p=0$) removes the pressure null space.
    \item Diagnostic option \code{Jexact}: given jumps enter the right-hand side instead of
      the reconstructed ones.
  \end{itemize}
\item \code{SolveIIM} (line~218) solves with a sparse direct solver and returns $z$,
  $\mathbf v=z/\eta$ and $p$.
\end{enumerate}
This first version builds the matrix from explicit triplets. Porting it to the StagFDTools
kernel style (AD-local residuals writing through \code{K[a][b]}) is left for when the method is
settled.

\paragraph{Test} (\code{test/test\_iim\_assembly.jl}). The test problem is a simple shear
parallel to an inclined $\Gamma$, with a velocity offset. Its velocity is piecewise linear, with
a kink and $[z]\neq0$. The IIM reproduces it to $10^{-9}$ (velocity) and $10^{-8}$ (pressure),
including where $\Gamma$ meets the domain boundary. Without corrections it does not.

\paragraph{Comparison.} The example \code{examples/\_Stokes\_IIM/IIM\_Step4\_Inclusion.jl}
compares three methods:
\begin{itemize}
  \item standard FD, i.e.\ the StagFDTools \code{Stokes} module with phase-ratio averaged
    viscosity, in \code{StandardFD\_Schmid2003.jl};
  \item the IIM;
  \item the IIM with exact jumps.
\end{itemize}
All three use the same exact Dirichlet velocities, $L=10$, $r_c=2$ and $\eta_{mat}=1$. Pressure
errors are computed after removing the mean of both the numerical and the exact pressure.

\begin{figure}[h]
  \centering
  \includegraphics[width=\textwidth]{figures/IIM_Step4_Inclusion.png}
  \caption{Viscous inclusion test.
  \textbf{Top:} $L_\infty$ and $L_1$ pressure errors against $h$ ($n_c=20$--320) for
  $\eta_{inc}/\eta_{mat}=10^{-3}$ (blue), $10^3$ (orange) and $10^6$ (green). Squares: standard
  FD; circles: IIM; triangles, dashed: IIM with exact jumps. The IIM curves for $10^3$ and $10^6$
  coincide, and so do the three exact-jump curves.
  \textbf{Bottom left:} pressure along $y=\Delta y/2$ at $n_c=80$, $\eta_{inc}/\eta_{mat}=10^3$.
  \textbf{Bottom right:} $\log_{10}|P-P^{ex}|$ for the same case.}
  \label{fig:step4}
\end{figure}

\begin{table}[h]
\centering
\small
\begin{tabular}{@{}r rr rr rr@{}}
\toprule
& \multicolumn{2}{c}{standard FD} & \multicolumn{2}{c}{IIM} & \multicolumn{2}{c}{IIM, exact jumps}\\
\cmidrule(lr){2-3}\cmidrule(lr){4-5}\cmidrule(l){6-7}
$n_c$ & $L_\infty$ & $L_1$ & $L_\infty$ & $L_1$ & $L_\infty$ & $L_1$\\
\midrule
 20 & 1.33 & 8.8e-2 & 1.13 & 1.3e-1 & 0.117  & 2.0e-2 \\
 40 & 1.72 & 9.1e-2 & 0.87 & 3.4e-2 & 0.174  & 1.1e-2 \\
 80 & 7.06 & 8.1e-2 & 0.70 & 2.7e-2 & 0.104  & 4.2e-3 \\
160 & 9.54 & 7.3e-2 & 0.40 & 7.9e-3 & 0.075  & 1.7e-3 \\
320 & 8.38 & 3.1e-2 & 0.27 & 5.1e-3 & 0.044  & 6.1e-4 \\
\bottomrule
\end{tabular}
\caption{Pressure errors for $\eta_{inc}/\eta_{mat}=10^3$. The IIM values for $10^6$ agree to
within 1\%. For $10^{-3}$, standard FD reaches $L_\infty=2.7$ and the IIM $0.19$ at $n_c=320$.}
\label{tab:step4}
\end{table}

\paragraph{Discussion.}
\begin{itemize}
  \item \emph{Standard FD} does not converge in $L_\infty$. Its error grows from 1.3 to about 8
    with resolution at high contrast; it is located in the cells cut by $\Gamma$, where the
    averaged viscosity smears the pressure jump over about two cells (Fig.~\ref{fig:step4},
    profile). $L_1$ decreases slowly.
  \item \emph{The IIM converges in $L_\infty$}, from 1.1 to 0.27 between $n_c=20$ and 320
    (mean order about 0.5, local orders 0.3--0.8). At $n_c=320$ the error is 30 times smaller
    than with standard FD. $L_1$ decreases at about $O(h)$, with scatter. The jump stays sharp
    within one cell.
  \item \emph{No degradation with contrast.} The IIM errors for $10^3$ and $10^6$ are identical,
    as anticipated in step~2: the error of the jumps does not depend on the contrast.
  \item \emph{Where the error comes from.} With exact jumps, the errors are 3--10 times smaller,
    but $L_\infty$ converges at a similar rate (about 0.5--0.8). Two causes act together:
    \begin{itemize}
      \item The first-order corrections leave an $O(1)$ truncation error at the irregular
        equations, because they omit $\tfrac12[\nabla^2z]\delta^2$ and $[\nabla p]\delta$. This
        limits the order.
      \item The reconstructed $\partial_t\mathbf v_\Gamma$ (step~2) multiplies the error
        constant. The remaining error sits inside the inclusion, next to $\Gamma$
        (Fig.~\ref{fig:step4}, error map).
    \end{itemize}
  \item \emph{Diagnostic detail.} The exact jumps are evaluated at $x_\Gamma$ projected onto the
    circle. Using $x_\Gamma$ itself, $O(h^4)$ off $\Gamma$, gives errors of up to 35 at $10^6$,
    because $\Delta\eta$ amplifies the kink of $\mathbf v$ (the effect seen in step~2).
\end{itemize}

%---------------------------------------------------------------
\section{Step 5: quadratic fits and second-order corrections}\label{sec:s5}

Step~4 showed two limitations of the first-order scheme. First, the corrections omit
$\tfrac12[\nabla^2z]\delta^2$ and $[\nabla p]\delta$, which bounds the order. Second, the linear
fits give $\partial_s\mathbf v_\Gamma$ only to $O(h)$, which inflates the error constant. Step~5
addresses both, as separate options so that their effects can be compared.

\paragraph{Quadratic fits (option 1).} On each side, \code{QuadraticFit}
(\code{Trace.jl}, line~60) fits the ~10 same-side nodes closest to $x_\Gamma\pm L\mathbf n$ with
the basis $[1,\delta x,\delta y,\delta x^2/2,\delta y^2/2,\delta x\,\delta y]$, centred at
$x_\Gamma$, by least squares. If the normal matrix is ill-conditioned, it adds nodes. The fit
gives:
\begin{itemize}
  \item $z_s$ to $O(h^3)$ and $\partial_sz_s=\mathbf t\cdot\nabla z_s$ to $O(h^2)$;
  \item the second derivative along the curved interface,
    $\partial_{ss}z_s=\mathbf t^\top(\nabla\nabla z_s)\mathbf t-\kappa\,\mathbf n\cdot\nabla z_s$,
    to $O(h)$.
\end{itemize}
As before, $\partial_{ss}\mathbf v_\Gamma=(\partial_{ss}z_-+\partial_{ss}z_+)/(\eta_-+\eta_+)$.
The curvature $\kappa=\nabla\cdot(\nabla\varphi/|\nabla\varphi|)$ is computed at every crossing
by central differences on the fine lattice (\code{κΦ}, \code{Geometry.jl}, line~102). It is
$O(h^2)$-accurate, as checked in \code{test\_iim\_geometry.jl}.

\paragraph{Second-order jumps (option 2).} \code{InterfaceJumps2} (\code{Jumps.jl}, line~47)
returns $[\nabla\nabla z]$ and $[\nabla p]$ in addition to the first-order jumps. Landry
removed his second-order version and left its curvature terms as a TODO, so these jumps were
derived here:
\begin{itemize}
  \item The first-order relations hold at every point of $\Gamma$. With $\mathbf w=\partial_s\mathbf v$,
    \[
    [\nabla z](s)=\Delta\eta\bigl(\mathbf w\,\mathbf t^\top-((\mathbf t\cdot\mathbf w)\mathbf n+(\mathbf n\cdot\mathbf w)\mathbf t)\,\mathbf n^\top\bigr).
    \]
  \item Differentiating along $\Gamma$, using $d\mathbf t/ds=-\kappa\mathbf n$ and
    $d\mathbf n/ds=\kappa\mathbf t$, gives all second derivatives that contain $\partial_t$:
    $[\partial_t\partial_\beta z_a]=d[\nabla z]_{a\beta}/ds$.
  \item The two remaining second derivatives, and $[\nabla p]$, follow from:
    \begin{itemize}
      \item $\partial_n(\nabla\cdot z)=0$, which gives $[\partial_{nn}z_n]=-[\partial_{nt}z_t]$;
      \item the tangential derivative of $[p]=-2\Delta\eta\,\mathbf t\cdot\mathbf w$, which
        gives $[\partial_tp]$;
      \item the momentum equations $\nabla^2z=\nabla p$, which give $[\partial_{nn}z_t]$ and
        $[\partial_np]$.
    \end{itemize}
  \item No body-force jump is assumed. Gravity with a density jump would add $[\rho]\mathbf g$
    to $[\nabla p]$.
\end{itemize}
\code{test\_iim\_jumps.jl} checks these jumps against the exact second derivatives of the Schmid
solution on both sides (ForwardDiff Hessians at $r_c(1\pm10^{-9})$). The errors are $10^{-8}$
to $10^{-12}$, for 3 contrasts, with and without simple shear.

\paragraph{Assembly.} \code{AssembleIIM(\dots; order=2, fit=:quadratic)} (\code{Assembly.jl},
line~110) extends each correction to
\[
q_k\;\to\;q_k+s\bigl([q]+[\nabla q]\cdot\delta+\tfrac12\delta^\top[\nabla\nabla q]\,\delta\bigr),
\qquad p_k\;\to\;p_k+s\bigl([p]+[\nabla p]\cdot\delta\bigr).
\]
\code{JumpFunctionals} (line~35) now handles clouds of any size: about 20 nodes per component
for quadratic fits. \code{test\_iim\_assembly.jl} adds a pressure-driven parallel flow along an
inclined $\Gamma$, an exact Stokes solution with a quadratic velocity. It is reproduced to
round-off at both orders. Its second-order jumps vanish, so it checks that the quadratic fits
and the new terms keep exactness; the second-order jumps themselves are tested against Schmid.

\paragraph{Results.} The example \code{examples/\_Stokes\_IIM/IIM\_Step5\_SecondOrder.jl}
compares four IIM variants with standard FD:
\begin{itemize}
  \item A: linear fits, order 1 (step~4);
  \item B: quadratic fits, order 1 (option 1);
  \item C: quadratic fits, order 2 (options 1 and 2);
  \item D: order 2 with exact jumps, which shows the accuracy of the corrections alone.
\end{itemize}
The shared driver is \code{IIM\_Schmid2003.jl}.

\begin{table}[h]
\centering
\small
\begin{tabular}{@{}r r rrrr r rr@{}}
\toprule
& FD & A & B & C & D & & C, $10^{-3}$ & C, $10^{6}$\\
\cmidrule(lr){2-6}\cmidrule(l){8-9}
$n_c$ & \multicolumn{5}{c}{$L_\infty(P)$, $\eta_{inc}/\eta_{mat}=10^3$} & ord.\ C & \multicolumn{2}{c}{$L_\infty(P)$}\\
\midrule
 20 & 1.33 & 1.13 & 0.584 & 0.569 & 0.055   & --   & 0.392  & 0.571 \\
 40 & 1.72 & 0.87 & 0.550 & 0.314 & 0.032   & 0.86 & 0.260  & 0.315 \\
 80 & 7.06 & 0.70 & 0.210 & 0.133 & 0.0096  & 1.24 & 0.108  & 0.134 \\
160 & 9.54 & 0.40 & 0.131 & 0.055 & 0.0027  & 1.27 & 0.054  & 0.055 \\
320 & 8.38 & 0.27 & 0.048 & 0.013 & 0.00068 & 2.05 & 0.0071 & 0.013 \\
\midrule
$L_1$, 320 & 3.1e-2 & 5.1e-3 & 6.0e-4 & 3.4e-4 & 5.4e-5 & & 4.1e-4 & 3.4e-4\\
\bottomrule
\end{tabular}
\caption{Pressure errors for the inclusion test. A: linear fits, order 1; B: quadratic fits,
order 1; C: quadratic fits, order 2; D: order 2, exact jumps. ``ord.\ C'' is the local
$L_\infty$ order of C at $10^3$.}
\label{tab:step5}
\end{table}

\begin{figure}[h]
  \centering
  \includegraphics[width=\textwidth]{figures/IIM_Step5_SecondOrder.png}
  \caption{Second-order IIM.
  \textbf{Top:} $L_\infty$ and $L_1$ pressure errors against $h$. Colours: standard FD (blue),
  A (orange), B (pink), C (green), D (grey, dashed). Markers: $\eta_{inc}/\eta_{mat}=10^3$
  (circles), $10^{-3}$ (triangles), $10^6$ (squares).
  \textbf{Bottom left:} pressure along $y=\Delta y/2$ at $n_c=80$, $\eta_{inc}/\eta_{mat}=10^3$.
  \textbf{Bottom right:} $\log_{10}|P-P^{ex}|$ for A and C.}
  \label{fig:step5}
\end{figure}

\begin{figure}[h]
  \centering
  \includegraphics[width=\textwidth]{figures/IIM_PressureFields.png}
  \caption{Pressure fields at $n_c=40$, 80 and 160 (rows), $\eta_{inc}/\eta_{mat}=10^3$, zoomed on
  the inclusion (script \code{IIM\_PressureFields.jl}). Columns: standard FD, IIM variants A, B and
  C, and the exact solution. All panels share the colour scale of the exact solution; lime and
  magenta mark cells more than 10\% beyond the exact range. Each title gives the $L_\infty$
  error.}
  \label{fig:fields}
\end{figure}

\paragraph{Discussion.}
\begin{itemize}
  \item \emph{Second-order IIM (C).} $L_\infty(P)$ decreases from 0.57 to 0.013 between
    $n_c=20$ and 320. That is a mean order of 1.35, and the local order reaches 2.05 at the
    finest step. $L_1$ converges at about $O(h^2)$. At $n_c=320$ the $L_\infty$ error is 630 times
    smaller than with standard FD and 20 times smaller than with the first-order IIM.
  \item \emph{Each option contributes.}
    \begin{itemize}
      \item Quadratic fits alone (B) reduce the error by a factor of 1.6--6 compared with A, but
        keep the first-order behaviour.
      \item Adding the second-order corrections (C) changes the order.
    \end{itemize}
  \item \emph{Limit of the corrections.} With exact jumps (D), $L_\infty(P)$ converges at
    $O(h^2)$ from $n_c=80$ on. The remaining gap between C and D (a factor of 10--20) comes
    from $\partial_{ss}\mathbf v_\Gamma$, which the quadratic fits give only to $O(h)$; it enters
    $[\nabla\nabla z]$ and $[\nabla p]$. Cubic fits would close part of this gap, at the cost of
    larger clouds.
  \item \emph{Contrast.} C gives the same errors for $10^3$ and $10^6$, and similar or slightly
    smaller ones for $10^{-3}$. The method does not degrade at high contrast.
  \item \emph{Pressure fields} (Fig.~\ref{fig:fields}). Standard FD shows spikes of both signs
    in the cells cut by $\Gamma$. They narrow but grow with resolution, which explains the
    non-converging $L_\infty$. All IIM variants keep the jump sharp without spikes. At
    $n_c=160$, variant C is visually indistinguishable from the exact solution.
  \item \emph{Cost.} The second-order rows are wider (about 20 nodes per component and crossing
    instead of 6). They remain confined to the band along $\Gamma$, so the extra cost scales
    as $O(n_c)$.
\end{itemize}

%---------------------------------------------------------------
\section{Step 6: compressible viscous flow}\label{sec:s6}

\paragraph{Equations.} As in the StagFDTools \code{Stokes} module (\code{compressible=true}),
\[
\boldsymbol\tau=2\eta\bigl(\dot{\boldsymbol\varepsilon}-\tfrac13(\nabla\cdot\mathbf v)\mathbf I\bigr),
\qquad p=-\xi\,\nabla\cdot\mathbf v,
\]
with a bulk viscosity $\xi$ that may also jump across $\Gamma$. With $z=\eta\mathbf v$ and
$k=\eta/\xi$ on each side ($k=0$: incompressible),
\[
\nabla^2z+\tfrac13\nabla(\nabla\cdot z)-\nabla p=0,\qquad \nabla\cdot z+k\,p=0.
\]
The pressure level is determined, so the Lagrange multiplier is dropped.
\code{AssembleIIM(\dots; ξ=(neg, pos))} (\code{Assembly.jl}, line~110) uses the factor
$\tfrac13$ instead of 1 in front of the grad-div terms and adds $k\,p$ to the continuity rows.

\paragraph{Jump conditions.} Without incompressibility, $\partial_nv_n$ no longer follows from
$\partial_tv_t$. Continuity of the normal traction $C=\mathbf n\cdot\boldsymbol\sigma\mathbf n$
gives instead, on each side,
\[
\partial_nz_n=\frac{kC-(1-\tfrac23k)\,\eta\,\mathbf t\cdot\mathbf w}{1+\tfrac43k},\qquad
\partial_nz_t=S-\eta\,\mathbf n\cdot\mathbf w,\qquad
p=-C+\tfrac43\partial_nz_n-\tfrac23\partial_tz_t,
\]
with $\mathbf w=\partial_s\mathbf v$ and $S=\mathbf t\cdot\boldsymbol\sigma\mathbf n$. The
jumps therefore depend on the traction $\mathbf T=\boldsymbol\sigma\mathbf n$ on $\Gamma$, in
addition to $\mathbf v_\Gamma$ and $\mathbf w$:
\begin{itemize}
  \item \code{SideFields} (\code{Jumps.jl}, line~82) computes $\nabla z$, $\nabla\nabla z$, $p$
    and $\nabla p$ of one side at $x_\Gamma$ from $(\mathbf w,\partial_{ss}\mathbf v,\mathbf T,
    d\mathbf T/ds)$. The second derivatives follow as in step~5, by differentiating along
    $\Gamma$ and closing with the side equations, which combine into $\nabla^2z=(1+k/3)\nabla p$:
    \[
    \partial_np=\frac{\partial_{tt}z_n-\partial_{nt}z_t}{1+\tfrac43k},\qquad
    \partial_{nn}z_n=-\partial_{nt}z_t-k\,\partial_np,\qquad
    \partial_{nn}z_t=(1+\tfrac k3)\,\partial_tp-\partial_{tt}z_t.
    \]
  \item \code{InterfaceJumpsC} (line~119) returns the differences of the two sides. For $k=0$
    on both sides, $C$ and $S$ cancel and the result equals \code{InterfaceJumps2} (tested).
    When $\eta/\xi$ is the same on both sides (equal Poisson ratios), $C$ cancels from the
    first-order jumps.
  \item \code{test\_iim\_jumps.jl} checks them against the exact Moutzouris (2026) solution
    of \code{ExactFieldSolutions}, for the circle and the ellipse and three $(\eta,\xi)$ pairs.
    First-order jumps agree to $10^{-7}$, second-order ones to better than $10^{-3}$ relative,
    the accuracy of the finite-difference reference.
\end{itemize}

\paragraph{Traction trace.} $\mathbf T$ and $d\mathbf T/ds$ come from the same one-sided fits:
\code{InterfaceTrace} now also returns one-sided gradient and Hessian weights, and
\code{PressureTrace} (\code{Trace.jl}, line~122) fits $p$ on each side. Each side gives
$\boldsymbol\sigma_s=-p_s\mathbf I+\nabla z_s+\nabla z_s^\top-\tfrac23(\nabla\cdot z_s)\mathbf I$,
and $\mathbf T$ is the average of the two sides. The interface clouds therefore include
pressure nodes. \code{test\_iim\_assembly.jl} adds two tests:
\begin{itemize}
  \item a uniform strain rate on each side of an inclined $\Gamma$, with different $\eta/\xi$,
    reproduced to round-off with linear and quadratic fits and at both orders;
  \item a regression on the Moutzouris circle at $n_c=40$, where $\nabla\cdot\mathbf v\neq0$ and
    curvature make the grad-div factor matter (the other tests are blind to it).
\end{itemize}

\paragraph{Test.} \code{examples/\_Stokes\_IIM/IIM\_Step6\_Compressible.jl} (drivers
\code{IIM\_Moutzouris2026.jl} and \code{StandardFD\_Moutzouris2026.jl}) solves the Moutzouris
inclusion in $[-0.5,0.5]^2$ with exact Dirichlet velocities: $\eta_{mat}=1$, $\xi_{mat}=2$,
pure shear $\dot\varepsilon=1$ plus far-field volumetric rate $\dot\zeta=0.5$. Two inclusions:
\begin{itemize}
  \item stiff: $\eta_{inc}=10^3$, $\xi_{inc}=10^5$ ($\eta/\xi$ differs from the matrix, so the
    traction trace matters);
  \item weak: $\eta_{inc}=10^{-3}$, $\xi_{inc}=5\cdot10^{-4}$.
\end{itemize}
Circle of radius 0.2, and ellipse with semi-axes $0.4\times0.2$ and the far field rotated by
$\pi/6$. The pressure is compared without removing any mean; the exact $|P|$ reaches 3.4--5.2.
Exact jumps (D) are evaluated at the projection of $x_\Gamma$ on the exact interface.

\begin{table}[h]
\centering
\small
\begin{tabular}{@{}r rrrr rr c rrrr rr@{}}
\toprule
& \multicolumn{6}{c}{circle} && \multicolumn{6}{c}{ellipse}\\
\cmidrule(lr){2-7}\cmidrule(l){9-14}
& \multicolumn{4}{c}{$10^3$} & \multicolumn{2}{c}{$10^{-3}$} && \multicolumn{4}{c}{$10^3$} & \multicolumn{2}{c}{$10^{-3}$}\\
$n_c$ & FD & A & C & D & FD & C && FD & A & C & D & FD & C\\
\midrule
 20 & 1.29 & 0.94 & 0.504 & 0.021   & 1.38 & 0.371  && 6.26 & 0.58 & 0.752 & 0.085  & 2.64 & 0.533\\
 40 & 1.65 & 0.77 & 0.260 & 0.016   & 2.66 & 0.203  && 6.08 & 0.78 & 0.416 & 0.029  & 3.01 & 0.232\\
 80 & 4.44 & 0.52 & 0.107 & 0.0045  & 3.35 & 0.071  && 6.57 & 0.46 & 0.242 & 0.0074 & 3.12 & 0.091\\
160 & 6.01 & 0.26 & 0.036 & 0.0010  & 3.51 & 0.038  && 6.55 & 0.38 & 0.088 & 0.0040 & 3.43 & 0.048\\
320 & 5.38 & 0.18 & 0.0098 & 0.00035 & 3.70 & 0.0059 && 8.04 & 0.28 & 0.037 & 0.0011 & 3.67 & 0.014\\
\bottomrule
\end{tabular}
\caption{$L_\infty(P)$ for the compressible inclusion, $\eta_{inc}/\eta_{mat}=10^3$ and
$10^{-3}$. A: linear fits, order 1; C: quadratic fits, order 2; D: order 2, exact jumps.}
\label{tab:step6}
\end{table}

\begin{figure}[h]
  \centering
  \includegraphics[width=\textwidth]{figures/IIM_Step6_Compressible.png}
  \caption{Compressible inclusion: $L_\infty(P)$ against $h$ for the circle (left) and the
  ellipse (right). Markers: $\eta_{inc}/\eta_{mat}=10^3$ (circles) and $10^{-3}$ (triangles).}
  \label{fig:step6}
\end{figure}

\begin{figure}[h]
  \centering
  \includegraphics[width=0.92\textwidth]{figures/IIM_Step6_EllipseFields.png}
  \caption{Pressure for the compressible elliptical inclusion, $\eta_{inc}/\eta_{mat}=10^3$, at
  $n_c=40$, 80 and 160 (rows). Columns: standard FD, IIM A and C, exact. Same colour
  conventions as Fig.~\ref{fig:fields}.}
  \label{fig:ellipse}
\end{figure}

\paragraph{Discussion.}
\begin{itemize}
  \item \emph{Standard FD} does not converge in $L_\infty(P)$, for either shape or contrast: the
    error stays at 1.3--8, comparable to the pressure itself, with spikes in the cut cells
    (Fig.~\ref{fig:ellipse}).
  \item \emph{Second-order IIM (C)} converges in all four cases. For the circle, $L_\infty(P)$
    drops from 0.50 to 0.0098 (mean order 1.4, local order 1.9 at the finest step); at
    $n_c=320$ it is 550 times smaller than FD. This is the same accuracy as the incompressible
    test (0.013, Table~\ref{tab:step5}): compressibility costs nothing.
  \item \emph{Ellipse.} Errors of C are 1.5--4 times larger than for the circle, and C converges
    at a mean order of 1.1 (0.75 to 0.037), still 220 times below FD at $n_c=320$. The ellipse
    tips have a curvature of 10 (twice that of the circle), so the $O(h)$ error of
    $\partial_{ss}\mathbf v_\Gamma$ and of $d\mathbf T/ds$ is larger there.
  \item \emph{Exact jumps (D)} converge at a mean order of 1.5 for both shapes, with local
    orders up to 2, so the corrected operator handles compressibility and the ellipse; the gap
    to C (a factor of 20--35 at $n_c\ge160$) again comes from the interface traces.
  \item \emph{First-order IIM (A)} improves only slowly (0.94 to 0.18 for the circle, 0.58 to
    0.28 for the ellipse), as in the incompressible case.
  \item \emph{Weak inclusion.} C behaves the same as for the stiff one (0.0059 and 0.014 at
    $n_c=320$).
\end{itemize}

%---------------------------------------------------------------
\section{Step 7: stress form with StagFDTools kernels}\label{sec:s7}

\paragraph{Motivation.} The z-form operator $\nabla^2z+c\,\nabla(\nabla\cdot z)-\nabla p$ only holds
for an isotropic, linear viscosity that is constant in each phase, and it never forms a stress.
The StagFDTools \code{Stokes} module instead interpolates $V_x$ to the $V_y$ nodes and $V_y$ to
the $V_x$ nodes, evaluates all components of $\nabla\mathbf v$ at centroids and vertices,
computes $\boldsymbol\tau=\mathcal D\,\dot{\boldsymbol\varepsilon}$ with a (tangent) operator
$\mathcal D$, and takes $\nabla\cdot\boldsymbol\sigma$ as residual. This is what makes
anisotropy, non-linear creep and plasticity easy. Step~7 builds the IIM on that structure; the
z-form is kept for comparison.

\paragraph{Formulation.} \code{AssembleIIMStress} (\code{AssemblyStress.jl}, line~149), unknowns
$\mathbf v$ and $p$:
\begin{itemize}
  \item \code{velocity\_gradient} (lines~52 and 59) returns $(\partial_xv_x,\partial_yv_y,\partial_yv_x,
    \partial_xv_y)$ at a centroid or a vertex exactly as the \code{Stokes} module: direct
    differences, and differences of the 4-point averages $\bar V_x$ (at $V_y$ nodes) and
    $\bar V_y$ (at $V_x$ nodes);
  \item \code{residual} (lines~75--87): $\partial_x\sigma_{xx}+\partial_y\tau_{xy}$ with
    $\sigma_{xx}$ at the two centroids and $\tau_{xy}$ at the two vertices of a $V_x$ node, and
    likewise in $y$; continuity $\nabla\cdot\mathbf v+p/\xi=0$. Here
    $\mathcal D=2\eta\mathbf I$ on the deviatoric $(\dot\varepsilon_{xx},\dot\varepsilon_{yy},
    \dot\varepsilon_{xy})$.
  \item \emph{One-sided stencils.} In the equation of a node on side $s$, every stencil value on
    the other side is replaced by the extension of the side-$s$ field, and all stresses of that
    equation use the material of side $s$. The kernel therefore only sees smooth fields: the
    full strain-rate tensor at each centroid and vertex is one-sided, with no averaging across
    $\Gamma$. The wider stencil (up to $4\times4$ $V_y$ nodes for $x$-momentum) reaches ghost
    nodes near the boundary; their crossing is the one closest to them
    (\code{member\_crossing}, line~28).
  \item Each corrected value is a linear form in the interface cloud, and the row coefficients
    are obtained by ForwardDiff of the local residual, as in the AD-local kernels of
    StagFDTools. A non-linear $\mathcal D$ can later enter the same residual.
\end{itemize}

\paragraph{Which field to extend.} With $\delta=x_k-x_\Gamma$, two extensions are exact up to the
Taylor remainder:
\begin{align*}
\texttt{extend=:v}:&\quad v_k\to v_k+s\bigl([\nabla v]\delta+\tfrac12\delta^\top[\nabla\nabla v]\delta\bigr),\\
\texttt{extend=:z}:&\quad v_k\to\bigl(\eta_kv_k+s\bigl([z]+[\nabla z]\delta+\tfrac12\delta^\top[\nabla\nabla z]\delta\bigr)\bigr)/\eta_s.
\end{align*}
\begin{itemize}
  \item The $v$-jumps (\code{InterfaceJumpsV}, \code{Jumps.jl}, line~133) depend on the shear
    traction even for incompressible flow ($\partial_nv_t=S/\eta-\mathbf n\cdot\mathbf w$ on
    each side), so both extensions use the traction trace of step~6.
  \item The remainder of \code{:v} is $O(\delta^3\nabla^3v)$ of the other side. A stiff-side
    equation multiplies it by $\eta_s$, so the error grows with the contrast: at
    $\eta_{inc}/\eta_{mat}=10^3$ the stiff-side momentum residuals of the exact Schmid
    solution are $\approx\eta_{inc}h$ (213 at $n_c=40$), even with exact jumps.
  \item The remainder of \code{:z} is $O(\delta^3\nabla^3z)$, a stress gradient, independent
    of the contrast. This is why Landry works with $z=\eta v$; the stress form keeps $v$ as
    unknown and only extends through $z$. \code{:z} is the default.
\end{itemize}

\paragraph{Solver scaling.} With $v$ as unknown, the stiff-side columns are $\sim10^{-3}$ and the
stiff-side rows $\sim10^{3}$; UMFPACK then pivots poorly and the LU fill-in is 2.4 times larger
than for the z-form at $n_c=160$ (solve 18.6~s instead of 3.6~s). \code{SolveIIMStress} (line~246)
scales velocity columns by $\eta$ and momentum rows by $1/\eta$ before the solve, which restores
the z-form cost (both about 55~s at $n_c=320$, where the direct solver needs about 10~GB).

\paragraph{Tests.} \code{test\_iim\_stress.jl}: the incompressible kink and the compressible
uniform-strain solutions of step~6 are reproduced to round-off with both extensions, linear and
quadratic fits and both orders; a Moutzouris-circle regression checks \code{:z} (0.27 at
$n_c=40$) against \code{:v} (2.0).

\paragraph{Results.} \code{examples/\_Stokes\_IIM/IIM\_Step7\_StressForm.jl}; the drivers of
steps~5 and~6 take \code{form=:z|:stress}. Contrast $10^3$ in all cases.

\begin{table}[h]
\centering
\small
\begin{tabular}{@{}l rr rr rr r@{}}
\toprule
& \multicolumn{2}{c}{A} & \multicolumn{2}{c}{C} & \multicolumn{2}{c}{D} & C, \code{:v}\\
\cmidrule(lr){2-3}\cmidrule(lr){4-5}\cmidrule(lr){6-7}\cmidrule(l){8-8}
$n_c=160$ & z & stress & z & stress & z & stress & stress\\
\midrule
Schmid             & 0.404 & 0.445 & 0.0550 & 0.0564 & 0.0027 & 0.0021 & 2.77\\
Moutzouris circle  & 0.260 & 0.254 & 0.0357 & 0.0365 & 0.0010 & 0.0012 & 1.50\\
Moutzouris ellipse & 0.380 & 0.399 & 0.0882 & 0.0903 & 0.0040 & 0.0040 & 8.88\\
\bottomrule
\end{tabular}
\caption{$L_\infty(P)$ at $n_c=160$ for the z-form and the stress form (A: linear fits, order~1;
C: quadratic fits, order~2; D: order~2, exact jumps), and the stress form with the
$v$-extension. At $n_c=320$ (Schmid, C): 0.0133 (z) and 0.0166 (stress).}
\label{tab:step7}
\end{table}

\begin{figure}[h]
  \centering
  \includegraphics[width=\textwidth]{figures/IIM_Step7_StressForm.png}
  \caption{$L_\infty(P)$ against $h$ for the z-form (solid, circles) and the stress form (dashed,
  crosses), and for the stress form with the $v$-extension (red, dotted).}
  \label{fig:step7}
\end{figure}

\begin{figure}[p]
  \centering
  \includegraphics[width=\textwidth]{figures/IIM_Step7_PressureFields_Schmid.png}\\[6pt]
  \includegraphics[width=\textwidth]{figures/IIM_Step7_PressureFields_Ellipse.png}
  \caption{Pressure fields of the two IIM forms (quadratic fits, order~2) at $n_c=40$, 80 and 160
  (rows), $\eta_{inc}/\eta_{mat}=10^3$, script \code{IIM\_Step7\_PressureFields.jl}.
  \textbf{Top:} Schmid circle (incompressible, mean pressure removed).
  \textbf{Bottom:} compressible Moutzouris ellipse. Columns: standard FD, z-form, stress form,
  exact (same colour conventions as Fig.~\ref{fig:fields}), and $\log_{10}|P_z-P_{stress}|$ with
  its maximum.}
  \label{fig:step7fields}
\end{figure}

\paragraph{Discussion.}
\begin{itemize}
  \item \emph{Pressure fields} (Fig.~\ref{fig:step7fields}). Both IIM forms keep the pressure
    jump sharp without the spikes of standard FD, and are visually indistinguishable from each
    other and from the exact solution from $n_c=80$ on. Their difference converges: its maximum
    drops from 0.16 to 0.025 (Schmid) and from 0.034 to 0.0084 (ellipse) between $n_c=40$ and
    160. It is concentrated in the stiff inclusion and along $\Gamma$, where the two forms
    use different stencils and crossings for the corrections; in the matrix it is below
    $10^{-3}$ at $n_c=160$. For Schmid the difference inside the inclusion is smooth (four
    lobes, mirroring the far-field pressure pattern), i.e.\ a small, slowly varying offset
    rather than local noise.
  \item With the $z$-extension, the stress form reproduces the z-form in all three tests:
    within 3\% for C and within 25\% for D at $n_c=160$, with the same convergence. The
    StagFDTools structure (velocity interpolation, stresses at centroids and vertices) is
    therefore compatible with the IIM, provided every stencil is made one-sided.
  \item At coarse resolution the stress form is somewhat less accurate (Schmid, C: 0.89 against
    0.57 at $n_c=20$), because its wider stencil reaches members up to two cells across
    $\Gamma$, where the Taylor extension is less accurate.
  \item The $v$-extension does not converge at this contrast: extending through $z=\eta v$ is
    essential, and for a general rheology the natural analogue is to scale by an effective
    viscosity of each side at $\Gamma$.
  \item Assembly by AD is about 4--5 times slower than the hand-written z-form stencils, but the
    direct solve dominates.
\end{itemize}

\subsection*{Pressure at vertices: one-sided interpolation}

\paragraph{Where it is needed.} The pressure is a centroid unknown. The \code{Stokes} module
also needs it at vertices: the tangent operator $\mathcal D$ has a pressure column
(\code{𝐷.v[i][3,4] * Pt\_v} in \code{SMomentum\_x\_Generic}), which is non-zero for
pressure-dependent rheologies (Drucker--Prager yield, dilatancy, caps). The vertex value is the
4-point average of the surrounding centroids, \code{Pt\_v = av(Pt)}, exactly as $V_x$ is averaged
to the $V_y$ nodes. For linear viscous rheology this column vanishes, so none of the results above
uses a vertex pressure.

\paragraph{The problem.} The pressure jumps across $\Gamma$ ($[p]=-2\Delta\eta\,\mathbf t\cdot
\mathbf w$ for viscous flow, $O(1)$). At a vertex whose 4 centroids are not all on its side, the
plain average mixes $p_-$ and $p_+$ and is wrong by a fraction of $[p]$, independently of $h$. A
yield function evaluated there would see a wrong pressure in every cut cell.

\paragraph{One-sided interpolation.} The rule of the stress form applies unchanged: for a vertex on
side $s$, each centroid $k$ on the other side is replaced by the extension of the side-$s$
pressure before averaging,
\[
P_v=\frac14\sum_{k=1}^{4}\begin{cases}p_k & \text{side}(k)=s,\\[2pt]
p_k+s\bigl([p]+[\nabla p]\cdot\delta_k\bigr) & \text{otherwise,}\end{cases}
\qquad \delta_k=x_k-x_\Gamma,
\]
with the jumps at the crossing closest to $x_k$ between the vertex and the centroid (as for the
stencil values), evaluated on the current solution. The four values are then samples of the smooth
side-$s$ pressure, and $P_v$ is accurate to $O(h^2)$ plus the error of the jumps. In the AD kernel
this only means fetching the 4 centroid pressures through the same corrected getter as the
other stencil values. \code{VertexPressure} (\code{AssemblyStress.jl}, line~249) implements it as
a post-processing step; \code{test\_iim\_stress.jl} checks that it is exact for the
compressible uniform-strain solution, while the plain average is not.

\paragraph{Test.} \code{IIM\_Step7\_VertexPressure.jl} evaluates both averages on the converged
stress-form solutions (quadratic fits, order~2, contrast $10^3$) and compares them with the exact
pressure at the vertices (Table~\ref{tab:vertexp}, Fig.~\ref{fig:vertexp}).

\begin{table}[h]
\centering
\small
\begin{tabular}{@{}r rrr c rrr@{}}
\toprule
& \multicolumn{3}{c}{Schmid circle} && \multicolumn{3}{c}{Moutzouris ellipse}\\
\cmidrule(lr){2-4}\cmidrule(l){6-8}
$n_c$ & plain & one-sided & centroids && plain & one-sided & centroids\\
\midrule
 40 & 2.39 & 0.357 & 0.394 && 1.94 & 0.397 & 0.436\\
 80 & 2.17 & 0.208 & 0.195 && 1.86 & 0.229 & 0.241\\
160 & 2.07 & 0.064 & 0.056 && 1.85 & 0.098 & 0.090\\
\bottomrule
\end{tabular}
\caption{Maximum pressure error at the vertices next to $\Gamma$ (not all 4 centroids on the
vertex's side), plain and one-sided averages, and $L_\infty(P)$ at the centroids for reference.}
\label{tab:vertexp}
\end{table}

\begin{figure}[h]
  \centering
  \includegraphics[width=\textwidth]{figures/IIM_Step7_VertexPressure.png}
  \caption{Vertex pressure. \textbf{Left:} maximum error at the vertices next to $\Gamma$
  against $h/L$, plain (dashed) and one-sided (solid) averages. \textbf{Middle and right:}
  $\log_{10}$ of the vertex pressure error for the ellipse at $n_c=80$, plain and one-sided.}
  \label{fig:vertexp}
\end{figure}

\paragraph{Discussion.} The plain average does not converge: its error stays at about 2 (half of
$|[p]|$) in a one-cell band along $\Gamma$, which is the pressure analogue of the cut-cell spikes
of standard FD. The one-sided average removes the band; its maximum error is that of the centroid
pressure (within 15\%) and converges with it. The same rule is needed for every quantity that is
averaged between centroids and vertices and jumps across $\Gamma$. The old stresses $\tau^0$ of
visco-elasticity are the next case: they are discontinuous across $\Gamma$ (only $\sigma\mathbf n$
is continuous), and their jumps are not given by the current jump conditions, so they are better
kept one-sided (stored per side, or recomputed from one-sided strain rates).

%---------------------------------------------------------------
\section{Step 8: matrix-free stress form with dynamic relaxation}\label{sec:s8}

The stress form of step~7 has also been solved matrix free, with dynamic relaxation and
Powell--Hestenes pressure iterations in the vectorised style of the standalone PHDR codes
(\code{standalone/PHDR\_codes/Stokes2D\_DYREL\_IIM\_Schmid.jl}). It reaches the accuracy of the
direct solvers ($L_\infty(P)=0.050$ at $n_c=160$, contrast $10^3$). The method, the choice of
penalty and the results are in a separate report, \code{docs/IIM\_DYREL/IIM\_DYREL.pdf}.

%---------------------------------------------------------------
\section{Step 9: free surface without sticky air}\label{sec:s9}

\paragraph{Test.} Gravity-driven relaxation of a Gaussian topography: incompressible fluid,
$\eta=\rho=1$, $\mathbf g=(0,-1)$, box $x\in[-3,3]$, $y\ge-5$, free surface
$y=h(x)=e^{-x^2/\sigma^2}$ with $\sigma=0.5$, instantaneous Stokes solution. The references are a
FEM solution on a body-fitted triangular mesh (\code{FCFV\_Geodynamics.jl}) and an FD solution on
a deformed (surface-fitted) staggered grid (\code{FullStaggeredGrid}). In both, the sides and the
bottom are no slip, so the IIM runs use no slip as well.

\paragraph{Method.} A regular MAC grid covers the box up to $y=1.5$; the surface cuts it. The
unknowns are the nodes in the fluid, $\varphi=y-h(x)\le0$, and the velocity faces next to a
fluid centroid (step~10 shows why), with the standard stencils
($\boldsymbol\sigma=-p\mathbf I+2\eta(\dot{\boldsymbol\varepsilon}-\frac13\nabla\cdot\mathbf v\,\mathbf I)$).
Every other stencil value above the surface is a ghost, the Taylor extension of the fluid solution from the closest point $\mathbf x_\Gamma$ of the surface,
\[
\mathbf v_k=\mathbf v_\Gamma+\nabla\mathbf v\,\delta+\tfrac12\delta^\top\nabla\nabla\mathbf v\,\delta,
\qquad p_k=p_\Gamma+\nabla p\cdot\delta,\qquad \delta=\mathbf x_k-\mathbf x_\Gamma.
\]
This is the one-sided half of the jump machinery:
\begin{itemize}
  \item $\mathbf v_\Gamma$, $\partial_s\mathbf v$ and $\partial_{ss}\mathbf v$ come from one-sided
    cubic least-squares fits of 24 fluid nodes (\code{QuadraticFit} with \code{degree=3}; quadratic
    fits are one order less accurate, step~10, and need 16 nodes, see the robustness paragraph);
  \item the normal derivatives, $p_\Gamma$ and $\nabla p_\Gamma$ follow from \code{SideFields}
    (step~6) with zero traction, $\mathbf T=d\mathbf T/ds=0$. For incompressible flow this gives
    $p_\Gamma=2\eta\,\partial_n v_n=-2\eta\,\partial_s v_s$;
  \item gravity enters the second-order terms: \code{SideFields} takes a body force $\mathbf f$,
    which adds $f_n/(1+\frac43k)$ to $\partial_n p$ and $-f_t$ to $\partial_{nn}z_t$. For a flat
    surface at rest this gives the exact hydrostatic ghost, $p_k=-\rho g\,\delta_y$.
  \item pressure ghosts are extrapolated linearly along the grid line of the stencil, through
    the fluid neighbour $p_k$ and $p_\Gamma$ at the crossing (fraction $s$ of the segment):
    $p_g=p_k+(p_\Gamma-p_k)/s$, as Gibou et al.\ (2002) do for Dirichlet conditions. This uses
    only first derivatives of the fits and ties $p_g$ to the fluid pressure (the Taylor ghost
    $p_\Gamma+\nabla p\cdot\delta$ is 10 times less accurate in step~10).
\end{itemize}
Each ghost value is a linear form in the fluid unknowns, so the residuals are assembled directly
and solved with a sparse direct solver. The traction condition fixes the pressure level: there
is no null space. Code: \code{examples/\allowbreak \_Stokes\_IIM/\allowbreak IIM\_FreeSurface.jl};
figures: \code{IIM\_FreeSurface\_Figures.jl}.

\paragraph{Results.} A flat surface between grid lines gives $\mathbf v=0$ and the hydrostatic
pressure to $10^{-12}$. Mass is conserved to machine precision ($\int v_y\,dx\approx10^{-15}$ at
every depth). Table~\ref{tab:fs} and Fig.~\ref{fig:fs}.

\begin{table}[h]
\centering
\small
\begin{tabular}{@{}r r ccc c ccc@{}}
\toprule
& & \multicolumn{3}{c}{no slip} && \multicolumn{3}{c}{free slip}\\
\cmidrule(lr){3-5}\cmidrule(l){7-9}
$n_x$ & unknowns & $v_y$(top) & $\max|v_x|$ & $\max v_y$ && $v_y$(top) & $\max|v_x|$ & $\max v_y$\\
\midrule
 48 &   5\,896 & $-0.2012$ & 0.0540 & 0.0644 && $-0.2293$ & 0.0532 & 0.0930\\
 96 &  23\,654 & $-0.2016$ & 0.0544 & 0.0654 && $-0.2296$ & 0.0535 & 0.0934\\
192 &  94\,750 & $-0.2016$ & 0.0545 & 0.0655 && $-0.2294$ & 0.0536 & 0.0935\\
384 & 379\,256 & $-0.2017$ & 0.0546 & 0.0655 && $-0.2294$ & 0.0537 & 0.0936\\
\bottomrule
\end{tabular}
\caption{Free-surface test: vertical velocity at the top of the bump and extrema of the cell
velocities. One direct solve, about 6~s at $n_x=384$.}
\label{tab:fs}
\end{table}

\begin{figure}[p]
  \centering
  \setlength{\tabcolsep}{2pt}
  \begin{tabular}{@{}cccc@{}}
  & FEM, body-fitted & FD, deformed grid & IIM, regular grid ($n_x=96$)\\
  \rotatebox{90}{\hspace{1.2cm}$v_x$} &
  \includegraphics[width=0.31\textwidth]{figures/FreeSurfaceTest_FEM_Vx.png} &
  \includegraphics[width=0.31\textwidth]{figures/FreeSurfaceTest_FDM_Vx.png} &
  \includegraphics[width=0.31\textwidth]{figures/FreeSurfaceTest_IIM_Vx.png}\\
  \rotatebox{90}{\hspace{1.2cm}$v_y$} &
  \includegraphics[width=0.31\textwidth]{figures/FreeSurfaceTest_FEM_Vy.png} &
  \includegraphics[width=0.31\textwidth]{figures/FreeSurfaceTest_FDM_Vy.png} &
  \includegraphics[width=0.31\textwidth]{figures/FreeSurfaceTest_IIM_Vy.png}\\
  \rotatebox{90}{\hspace{1.2cm}$p$} &
  \includegraphics[width=0.31\textwidth]{figures/FreeSurfaceTest_FEM_P.png} &
  \includegraphics[width=0.31\textwidth]{figures/FreeSurfaceTest_FDM_P.png} &
  \includegraphics[width=0.31\textwidth]{figures/FreeSurfaceTest_IIM_P.png}
  \end{tabular}
  \caption{Free-surface test, no slip on the sides and the bottom; same colour ranges in each
  row ($v_x\in[-0.055,0.055]$, $v_y\in[-0.2,0.1]$, $p\in[0,5]$). IIM: cells clipped at the
  surface; in cut cells whose centroid is in the void, the ghost (extended) values are shown.}
  \label{fig:fs}
\end{figure}

\paragraph{Discussion.}
\begin{itemize}
  \item \emph{Agreement.} The IIM fields match the FEM solution: $v_y\approx-0.20$ at the top of
    the bump, uplift maxima of 0.065 near $x=\pm2$, $|v_x|\approx0.054$ at
    $(\pm0.9,-0.4)$, and nearly flat isobars. The deformed-grid FD solution shows the same
    velocities, but its isobars bulge under the bump. That script has a bulk viscosity
    $\eta_b=K\Delta t=1$, so its material is compressible, which may account for the difference.
  \item \emph{Boundary conditions.} With free slip on the sides, the uplift is largest at the
    walls (0.09) and the bump sinks faster ($-0.229$). Both references are no slip.
  \item \emph{Convergence.} $v_y$(top) varies by $5\cdot10^{-4}$ over $n_x=48$--384 and by
    $10^{-4}$ from 192 to 384. The largest pressure difference between these grids is
    $4\cdot10^{-3}$ (free slip) next to the surface and $2$--$5\cdot10^{-4}$ below $y=-0.5$; with
    no slip it is $5\cdot10^{-2}$ at the corners where the surface meets the walls (a singularity
    of the no-slip problem).
  \item \emph{Negative pressure.} The pressure is positive everywhere in the fluid: its minimum,
    0.02, is at the top of the bump. Along the flanks the surface pressure
    $p_\Gamma=2\eta\,\partial_nv_n$ is about 0.25 (Fig.~\ref{fig:fsps}). In the field figures, cut
    cells whose centroid is in the void show the extended pressure, slightly negative
    (a few hundredths at most) since it extrapolates the hydrostatic gradient above the surface.
  \item \emph{Stress-form kernel.} With Vx interpolated to the Vy nodes and Vy to the Vx nodes, as
    in the StagFDTools stress form, the solution differs by $3\cdot10^{-11}$: for isotropic
    viscosity the interpolated components enter $\boldsymbol\tau$ with zero weight. They add 44\%
    ghost nodes and explicit zeros to the matrix, and matter once $\mathcal D$ couples the
    components (anisotropy, plasticity).
  \item \emph{Resolution.} The radius of curvature at the top of the bump is
    $\sigma^2/2=0.125$, one cell at $n_x=48$; with cubic fits $v_y$(top) is nevertheless within
    $5\cdot10^{-4}$ of the finest grid there (quadratic fits: $2\cdot10^{-3}$).
  \item \emph{Next.} Time stepping (advecting $\varphi$, initialising nodes that leave the void
    with the same extension) and the implicit surface correction (FSSA) against the drunken-sailor
    instability. That correction is a Robin condition
    $\mathbf T=-\Delta t\,(\rho\mathbf g\cdot\mathbf n)(\mathbf v\cdot\mathbf n)\,\mathbf n$, which
    \code{SideFields} accepts through $\mathbf T$.
\end{itemize}

\paragraph{Robustness: defects found on the way} (more in step~10).
\begin{itemize}
  \item \emph{Pressure null space.} The normal-stress condition reaches the solution only
    through the pressure ghosts, i.e.\ through momentum equations whose stencil contains a void
    centroid. With a flat surface lying exactly on a row of Vy nodes, those nodes were ghosts and no
    equation reached the void: the pressure level was free (error 4.3 in the hydrostatic test).
    Nodes on the surface ($\varphi=0$) are now fluid nodes. Making nodes within $\theta\Delta$ of
    the surface ghosts (a common small-cut remedy) recreates the problem and must not be used.
  \item \emph{Fit size.} With 10-point quadratic fits (6 coefficients), the surface pressure had
    non-converging spikes at $x\approx\pm0.63$, where the slope is $45^\circ$, at $n_x=96$ and
    384 (left of Fig.~\ref{fig:fsps}). The fits pass the conditioning check there but are close
    to degenerate on the staircase of nodes; the extension carries the noise into the solution.
    16-point fits remove the spikes. On a manufactured solution (smooth flow with the matching
    traction and body force, free slip) both converge at second order, the 16-point fits with
    twice larger errors ($|\mathbf v|$ error $1.7\cdot10^{-5}$, $p$ error $1.5\cdot10^{-4}$ at
    $n_x=384$). The cubic fits now used have 10 coefficients and take 24 nodes.
\end{itemize}

\begin{figure}[h]
  \centering
  \includegraphics[width=\textwidth]{figures/IIM_FreeSurface_SurfacePressure.png}
  \caption{Surface pressure $p_\Gamma=2\eta\,\partial_nv_n$ with 10- and 16-point quadratic fits, and the
  surface velocity, no slip. The rise at $x=\pm3$ is the corner between the surface and the
  no-slip walls.}
  \label{fig:fsps}
\end{figure}

%---------------------------------------------------------------
\section{Step 10: free-surface accuracy test with an exact solution}\label{sec:s10}

\paragraph{Test.} The test of Duretz et al.\ (2016, GJI 204, Fig.~3), also used in GMD 19,
8839 (2026, Figs~7--9): the inviscid circular inclusion of Schmid \& Podladchikov (2003),
radius 0.5 at the origin, seen through the box $[-0.25,0.25]\times[-0.75,-0.25]$. The matrix
($\eta=1$) fills the bottom of the box and its top is a curved free surface (zero traction is
exact for an inviscid inclusion). Exact velocities on the sides and the bottom, no gravity. Both
papers report the pressure error in L1 only. Neither states the far field; I used pure shear,
$\dot\varepsilon=-1$. Errors are taken at the centroids (pressure) and faces (velocity) in the
matrix: L1 $=\frac1V\int|e|\,dV$ (mean over the cells), L$_\infty=\max|e|$.

\paragraph{Variants.}
\begin{itemize}
  \item \emph{IIM free surface:} step~9 with a circular level set (exact closest point and
    curvature) and Dirichlet sides; cubic (default) or quadratic fits.
  \item \emph{Sticky air:} the inclusion filled with viscosity $\eta_{air}$, exact inviscid
    velocities on all sides (also on the top, in the air), compared with the inviscid solution
    after removing the mean pressure difference over the matrix (the pressure level is free).
    (a) Two-phase IIM, z-form of step~5 (\code{AssembleIIM}, order 2, cubic or quadratic fits),
    up to $n=160$ (memory); (b) standard FD with phase-averaged viscosity
    (\code{StandardFD\_Schmid2003.jl}).
\end{itemize}
Code: \code{examples/\allowbreak \_Stokes\_IIM/\allowbreak IIM\_FreeSurface\_Schmid.jl}.

\paragraph{Defects found and fixed.}
\begin{itemize}
  \item \emph{Surface pressure missing in columns where $\Gamma$ is nearly parallel to a grid
    line.} At the bottom of the circle, a row of centroids lies a fraction of a cell below
    $\Gamma$. The face above each of them is in the void, so it was a ghost without an
    equation, and $p_\Gamma$ (which enters only through the pressure ghosts of momentum
    equations) never reached these columns: their pressure was integrated horizontally from the
    ends of the row. Shifting the grid by $0.37\Delta$ raised L$_\infty(p)$ from 0.03 to 1.9 at
    $n=160$. Fix: every face next to a fluid centroid is an unknown with its momentum equation,
    as in ghost-fluid free-surface methods; its void neighbour centroid carries the pressure
    ghost. The error is then independent of the grid shift. This also improves step~9
    (manufactured solution: L$_\infty(p)$ $4.4\cdot10^{-5}$ instead of $1.5\cdot10^{-4}$ at
    $n_x=384$, cleanly second order).
  \item \emph{Dirichlet values in the void.} Where $\Gamma$ meets a side, wall nodes and mirror
    ghosts above $\Gamma$ took the exact solution at the wall, i.e.\ the inclusion's velocity
    rather than the extension of the matrix flow. They are now extensions like any void node.
  \item \emph{Dirichlet wall ghosts.} The tangential ghost across a Dirichlet wall was the linear
    mirror $2v_w-v_1$, which gave a first-order pressure error at the bottom corners of the box
    (visible in the error map of FD sticky air, which uses the same mirror). It is now quadratic,
    $(8v_w-6v_1+v_2)/3$.
  \item \emph{z-form near the boundary.} For a ghost node across $\Gamma$ next to the boundary,
    the crossing between it and its row node can lie outside the domain, where no trace is
    built; \code{pair\_crossing} now falls back to the nearest crossing (as the stress form did).
\end{itemize}

\paragraph{Error budget of the free surface.} To find what limits the accuracy, the fitted
quantities were replaced one at a time by exact ones (diagnostic options \code{vexact}: exact
velocity ghosts from the analytic continuation of the matrix flow; \code{pΓexact}: exact
$p_\Gamma$). With quadratic fits, at $n=160$:
\begin{center}
\small
\begin{tabular}{@{}l cc l@{}}
\toprule
& $p$ L1 & $p$ L$_\infty$ & behaviour\\
\midrule
quadratic fits                          & 2.8e-3 & 1.0e-2 & L$_\infty$ order 1.4--2, irregular\\
exact velocity ghosts                   & 3.7e-3 & 4.6e-3 & L$_\infty$ clean second order\\
exact $p_\Gamma$                        & 2.4e-3 & 6.8e-3 & uniform offset changes sign\\
both exact, linear wall ghosts          & 1.8e-4 & 3.0e-3 & L$_\infty$ first order, bottom corners\\
both exact, quadratic wall ghosts       & 3.3e-4 & 4.3e-4 & \\
cubic fits (no exact data)              & 1.2e-4 & 6.8e-4 & L1 third, L$_\infty$ second order\\
\bottomrule
\end{tabular}
\end{center}
The fitted velocity ghosts and the fitted $p_\Gamma$ contribute errors of the same size and
opposite sign to a uniform pressure offset; the velocity ghosts also make the L$_\infty$ rate
irregular. Neither the pressure ghost nor its order is the limit: quadratic pressure ghosts
(through two fluid centroids and $p_\Gamma$, Gibou \& Fedkiw 2005) change nothing. Cubic fits,
one order more accurate for $\partial_s\mathbf v$ and $\partial_{ss}\mathbf v$, reduce both
contributions: with them, L1$(p)$ is 150 times smaller at $n=640$.

\begin{table}[h]
\centering
\small
\begin{tabular}{@{}r cc cc c cc cc c cc@{}}
\toprule
& \multicolumn{4}{c}{IIM free surface} && \multicolumn{4}{c}{IIM sticky, $\eta_{air}=10^{-6}$} && \multicolumn{2}{c}{FD sticky, $10^{-4}$}\\
\cmidrule(lr){2-5}\cmidrule(lr){7-10}\cmidrule(l){12-13}
$n$ & $p$ L1 & $p$ L$_\infty$ & $\mathbf v$ L1 & $\mathbf v$ L$_\infty$ && $p$ L1 & $p$ L$_\infty$ & $\mathbf v$ L1 & $\mathbf v$ L$_\infty$ && $p$ L1 & $p$ L$_\infty$\\
\midrule
 20 & 5.7e-2 & 1.1e-1 & 3.2e-4 & 2.6e-3 && 2.5e-3 & 2.4e-2 & 7.5e-5 & 4.7e-4 && 1.8e-1 & 2.4\\
 40 & 5.8e-3 & 1.5e-2 & 2.2e-5 & 1.3e-4 && 3.9e-4 & 5.8e-3 & 2.0e-5 & 8.4e-5 && 1.1e-1 & 2.6\\
 80 & 8.5e-4 & 2.6e-3 & 4.3e-6 & 1.8e-5 && 7.5e-5 & 7.5e-4 & 3.8e-6 & 1.6e-5 && 5.9e-2 & 2.6\\
160 & 1.2e-4 & 6.8e-4 & 1.1e-6 & 3.6e-6 && 1.8e-5 & 3.3e-4 & 8.6e-7 & 4.0e-6 && 3.2e-2 & 2.7\\
320 & 1.3e-5 & 1.7e-4 & 2.4e-7 & 8.4e-7 && & & & && 1.7e-2 & 2.8\\
640 & 1.2e-6 & 3.6e-5 & 5.8e-8 & 2.0e-7 && & & & && 8.5e-3 & 2.8\\
\midrule
order & 2.8--3.5 & 2.0--2.2 & 2.0--2.1 & 2.0--2.4 && 2.1--2.7 & 1.2--2.9 & 2.2--2.4 & 2.0--2.5 && 0.9 & --\\
\bottomrule
\end{tabular}
\caption{Inviscid inclusion, cubic fits: errors against the exact solution. IIM free surface with
linear pressure ghosts (no pressure shift); sticky air with the mean pressure difference
removed. Orders from $n=80$ on for the free surface. With quadratic fits: free surface
$1.8\cdot10^{-4}$ / $9.3\cdot10^{-4}$ (L1 / L$_\infty$ of $p$ at $n=640$, second order); IIM sticky
$4.9\cdot10^{-5}$ / $2.0\cdot10^{-3}$ at $n=160$.}
\label{tab:fs2}
\end{table}

\begin{figure}[h]
  \centering
  \includegraphics[width=\textwidth]{figures/IIM_FreeSurface_Schmid.png}
  \caption{Inviscid inclusion. \textbf{Top:} pressure errors in L1 and L$_\infty$, velocity
  errors, against $h$. Black: IIM free surface (cubic, quadratic fits); blue: IIM sticky air
  (solid cubic, dashed quadratic fits); orange: FD sticky air (the two contrasts overlap). In the
  velocity panel, grey: free surface with quadratic fits. \textbf{Bottom:} $\log_{10}$ of the
  pressure error at $n=160$, cubic fits.}
  \label{fig:fs2}
\end{figure}

\paragraph{Discussion.}
\begin{itemize}
  \item \emph{Free surface.} With cubic fits, pressure converges at third order in L1 and second
    order in L$_\infty$ ($3.6\cdot10^{-5}$ at $n=640$), velocity at second order in both norms.
    Duretz et al.\ (2016) and the GMD study get first order in L1 and do not report L$_\infty$.
    The fits matter more than anything else: at $n=640$ quadratic fits give 150 times larger L1
    and 26 times larger L$_\infty$. Pressure ghosts: linear ghosts are 10 times more accurate in
    L1 than Taylor ghosts ($p_\Gamma+\nabla p\cdot\delta$), which also need $\nabla p_\Gamma$.
  \item \emph{FD sticky air.} First order in L1, no convergence in L$_\infty$ (2.7 in the cut
    cells), the same for $\eta_{air}=10^{-2}$ and $10^{-4}$: the error of the averaged viscosity
    dominates the modelling error of the air. This matches the papers.
  \item \emph{IIM sticky air.} With $\eta_{air}=10^{-2}$ the error stalls at L1 $=4.7\cdot10^{-3}$
    (L$_\infty=7\cdot10^{-2}$): the modelling error of a viscous air, not a discretisation error.
    With $10^{-4}$ it stalls from $n=160$ on (L1 $5.5\cdot10^{-5}$); with $10^{-6}$ it converges
    at second order. The IIM removes the discretisation error of sticky air and leaves only the
    modelling error, proportional to $\eta_{air}$. The price is the air unknowns and a much more
    expensive direct solve (18~s against 0.6~s at $n=160$).
  \item \emph{Why sticky air seemed to converge faster in L$_\infty$} (quadratic fits). Its
    maximum error sits in cells whose centroid lies almost on $\Gamma$ (0.01--0.03$\Delta$); it
    starts large and decreases quickly, a pre-asymptotic rate. A few cells away from $\Gamma$ its
    error converged at first order only. The free-surface error was instead a smooth error over
    the whole domain, set by the fits (error budget above), whose maximum moved between the bottom
    of the circle and the corners with the walls, hence the irregular rate. With cubic fits both
    are more accurate and comparable at $n=160$ (L$_\infty$ $6.8\cdot10^{-4}$ free surface,
    $3.3\cdot10^{-4}$ sticky air).
  \item \emph{Remaining limit of IIM sticky air.} At $n=160$ its largest error is where $\Gamma$
    meets the side walls (L$_\infty$ rate 1.2 between $n=80$ and 160). There the cubic fits fall
    back to quadratic ones (too few nodes) and the corrections use the nearest crossing inside
    the domain. Away from the walls the error is 2--4 times smaller.
\end{itemize}

%---------------------------------------------------------------
\section{Next steps}

\begin{itemize}
  \item \emph{Accuracy:} cubic fits for $\partial_{ss}\mathbf v_\Gamma$ and $d\mathbf T/ds$, to
    approach the exact-jump curve D (the gap is largest for the ellipse); density jumps
    ($[\rho]\mathbf g$ in $[\nabla p]$).
  \item \emph{Compressibility:} the out-of-plane strain rate $\dot\varepsilon_{zz}$ of the
    Moutzouris solution, and elastic compressibility $\beta(P-P^0)/\Delta t$, which adds a source
    $[\beta P^0/\Delta t]$ to the continuity jump.
  \item \emph{Robustness:} more geometries (SolCx with an inclined interface, sub-cell offsets
    of the inclusion, thin layers and their fallback).
  \item \emph{General rheology in the stress form:} replace the closed-form \code{SideFields} by
    a side solve for the missing normal derivatives (continuity of $\mathbf v$ and of
    $\boldsymbol\sigma\mathbf n$ with the tangent $\mathcal D$ of each side); a manufactured
    anisotropic test; power-law creep in the inclusion with Newton iterations.
  \item \emph{StagFDTools integration:} AD-local kernels writing through \code{K[a][b]}, GCR
    with the averaged-viscosity preconditioner, markers and moving interfaces.
\end{itemize}

\end{document}
